\section{嘉靖至隆庆的补充编年节点（一）}
这一组条目把此前尚未设导航入口的年度材料接回本章脉络。它们不是把同年材料机械等同为因果，而是在可核查的时间位置上，分别保留行动、行政记录与社会经验之间的距离。

\subsection{1529年前后的材料线索}

\eventanchor{ML-1529-001}{明·1529（嘉靖八年）}
\paragraph{1529（嘉靖八年）：嘉靖倭寇袭扰：温州数名将领因与倭寇勾结被流放} 这一节点应放回本章已经展开的权力、财政、地方社会与边地关系中理解。账本把它出现的条件概括为：前期边防、地方武装或跨区域资源调动的压力推动了该行动。 随后可追踪的变化是：它改变了后续调兵、补给或地方秩序的条件；战果和伤亡须分开核对。 可由《世宗实录》嘉靖八年条；《明史·兵志》及相关列传；边镇、地方志或对方材料 互相核对；不过，译写候选以编年材料定位；实录记载反映朝廷选择，崇祯期须另以奏疏、地方及敌对政权材料交叉。 因而这里标示的是材料能够支持的行动与制度位置，而不是对动机、地方执行或普通人经验的无保留复原。

\eventanchor{ML-1529-002}{明·1529（嘉靖八年）}
\paragraph{1529（嘉靖八年）：看到了一颗不吉利的彗星} 这一节点应放回本章已经展开的权力、财政、地方社会与边地关系中理解。账本把它出现的条件概括为：继承、官僚协调或皇权运作的压力使问题进入朝廷程序。 随后可追踪的变化是：它影响后续人事与政策议程；动机和派系标签不能由单一叙事裁定。 可由《世宗实录》嘉靖八年条；《明史·本纪》及相关列传；诏令、奏疏或会典版本 互相核对；不过，译写候选以编年材料定位；实录记载反映朝廷选择，崇祯期须另以奏疏、地方及敌对政权材料交叉。 因而这里标示的是材料能够支持的行动与制度位置，而不是对动机、地方执行或普通人经验的无保留复原。

\eventanchor{ML-1529-003}{明·1529（嘉靖八年）}
\paragraph{1529（嘉靖八年）：嘉靖朝礼制、内阁与人事的本年诏令和奏议形成中枢协调线索，留中与施行之间应予区分。} 这一节点应放回本章已经展开的权力、财政、地方社会与边地关系中理解。账本把它出现的条件概括为：继承、官僚协调或皇权运作的压力使问题进入朝廷程序。 随后可追踪的变化是：它影响后续人事与政策议程；动机和派系标签不能由单一叙事裁定。 可由《世宗实录》嘉靖八年条；《明史·本纪》及相关列传；诏令、奏疏或会典版本 互相核对；不过，桥接条目只标示可回查的年度问题与材料链；实录有中央选择性，崇祯期尤其需要奏疏、地方、朝鲜及满文材料交叉。 因而这里标示的是材料能够支持的行动与制度位置，而不是对动机、地方执行或普通人经验的无保留复原。

\eventanchor{ML-1529-004}{明·1529（嘉靖八年）}
\paragraph{1529（嘉靖八年）：祀典、学校、出版或讲学的本年材料记录文化制度活动，规范话语不等于社会普遍认同。} 这一节点应放回本章已经展开的权力、财政、地方社会与边地关系中理解。账本把它出现的条件概括为：制度、教育或知识传播的需要推动了相关行动或文本形成。 随后可追踪的变化是：它提供制度和传播史的锚点，不能据此推断社会整体接受。 可由《世宗实录》嘉靖八年条；《明史·选举志》或《艺文志》；文集、碑刻或书目材料 互相核对；不过，桥接条目只标示可回查的年度问题与材料链；实录有中央选择性，崇祯期尤其需要奏疏、地方、朝鲜及满文材料交叉。 因而这里标示的是材料能够支持的行动与制度位置，而不是对动机、地方执行或普通人经验的无保留复原。

\eventanchor{MM-1529-WANG-YANGMING-DEATH}{明·1529（嘉靖七年）}
\paragraph{1529（嘉靖七年）：王守仁自广西返乡途中病逝，其学术与军政文字随后由门人整理传播} 这一节点应放回本章已经展开的权力、财政、地方社会与边地关系中理解。账本把它出现的条件概括为：广西经略消耗与长期病疾限制了王守仁返乡后的活动，他的弟子网络已在江南和东南形成抄传、讲学与仕宦联系。 随后可追踪的变化是：其思想在书院、家塾和官场继续扩散并引发争论；这种传播不应被简化为全国士人立刻接受同一学说。 可由《明史·王守仁传》；《王阳明年谱》嘉靖七年条；《王阳明全集》序跋与遗书 互相核对；不过，死亡时间与行程大体可由年谱、文集和传记互校；门人编订文本有选择与后出修订，不能当作其每一行动的同期记录。 因而这里标示的是材料能够支持的行动与制度位置，而不是对动机、地方执行或普通人经验的无保留复原。

\eventanchor{ML-1530-001}{明·1530（嘉靖九年）}
\paragraph{1530（嘉靖九年）：赋役折色、盐课、河工或田土核查的本年文书为财政改革史提供节点，制度名称不可跨区套用。} 这一节点应放回本章已经展开的权力、财政、地方社会与边地关系中理解。账本把它出现的条件概括为：财政征解、交通工程或货币支付的压力使相关措施进入政务。 随后可追踪的变化是：其法定安排不等于地方执行，后续须以地域材料追踪负担与成效。 可由《世宗实录》嘉靖九年条；《明会典》相关条目；地方志、黄册/契约/河工材料 互相核对；不过，桥接条目只标示可回查的年度问题与材料链；实录有中央选择性，崇祯期尤其需要奏疏、地方、朝鲜及满文材料交叉。 因而这里标示的是材料能够支持的行动与制度位置，而不是对动机、地方执行或普通人经验的无保留复原。

\eventanchor{ML-1530-002}{明·1530（嘉靖九年）}
\paragraph{1530（嘉靖九年）：嘉靖朝礼制、内阁与人事的本年诏令和奏议形成中枢协调线索，留中与施行之间应予区分。} 这一节点应放回本章已经展开的权力、财政、地方社会与边地关系中理解。账本把它出现的条件概括为：继承、官僚协调或皇权运作的压力使问题进入朝廷程序。 随后可追踪的变化是：它影响后续人事与政策议程；动机和派系标签不能由单一叙事裁定。 可由《世宗实录》嘉靖九年条；《明史·本纪》及相关列传；诏令、奏疏或会典版本 互相核对；不过，桥接条目只标示可回查的年度问题与材料链；实录有中央选择性，崇祯期尤其需要奏疏、地方、朝鲜及满文材料交叉。 因而这里标示的是材料能够支持的行动与制度位置，而不是对动机、地方执行或普通人经验的无保留复原。

\eventanchor{ML-1530-003}{明·1530（嘉靖九年）}
\paragraph{1530（嘉靖九年）：北边警报、边镇防务和西南处置的本年军政材料标记边疆压力，敌我称谓与军数应回到原文。} 这一节点应放回本章已经展开的权力、财政、地方社会与边地关系中理解。账本把它出现的条件概括为：前期边防、地方武装或跨区域资源调动的压力推动了该行动。 随后可追踪的变化是：它改变了后续调兵、补给或地方秩序的条件；战果和伤亡须分开核对。 可由《世宗实录》嘉靖九年条；《明史·兵志》及相关列传；边镇、地方志或对方材料 互相核对；不过，桥接条目只标示可回查的年度问题与材料链；实录有中央选择性，崇祯期尤其需要奏疏、地方、朝鲜及满文材料交叉。 因而这里标示的是材料能够支持的行动与制度位置，而不是对动机、地方执行或普通人经验的无保留复原。

\eventanchor{ML-1530-004}{明·1530（嘉靖九年）}
\paragraph{1530（嘉靖九年）：灾异、赈济及地方工程的本年报告可与州县材料互校，救济命令不自动证明救济到达。} 这一节点应放回本章已经展开的权力、财政、地方社会与边地关系中理解。账本把它出现的条件概括为：自然风险与地方报告促成朝廷或地方的应对记录。 随后可追踪的变化是：赈济、迁徙和损失的实际范围仍须以受灾地区材料分别核验。 可由《世宗实录》嘉靖九年条；《明史·五行志》；受灾府县地方志与奏疏 互相核对；不过，桥接条目只标示可回查的年度问题与材料链；实录有中央选择性，崇祯期尤其需要奏疏、地方、朝鲜及满文材料交叉。 因而这里标示的是材料能够支持的行动与制度位置，而不是对动机、地方执行或普通人经验的无保留复原。

\eventanchor{MM-1530-SUBURBAN-RITES}{明·1530（嘉靖九年）}
\paragraph{1530（嘉靖九年）：嘉靖调整郊祀，分别祭天与祭地，并改建相应坛制} 这一节点应放回本章已经展开的权力、财政、地方社会与边地关系中理解。账本把它出现的条件概括为：大礼议后，嘉靖持续以礼制重申皇帝对宗法和祭祀解释权的裁断。 随后可追踪的变化是：北京国家祭祀空间随之重组，礼制改革也为嘉靖后期频繁斋醮和宫廷宗教活动奠定制度环境。 可由《世宗实录》嘉靖九年礼部条；《明史·礼志》；北京天坛、地坛营建与碑刻资料 互相核对；不过，礼制调整和工程命令可靠；仪式文本体现国家规范，不能说明社会普遍信仰或实际参与。 因而这里标示的是材料能够支持的行动与制度位置，而不是对动机、地方执行或普通人经验的无保留复原。

\eventanchor{ML-1531-001}{明·1531（嘉靖十年）}
\paragraph{1531（嘉靖十年）：大同遭到蒙古人的袭击} 这一节点应放回本章已经展开的权力、财政、地方社会与边地关系中理解。账本把它出现的条件概括为：前期边防、地方武装或跨区域资源调动的压力推动了该行动。 随后可追踪的变化是：它改变了后续调兵、补给或地方秩序的条件；战果和伤亡须分开核对。 可由《世宗实录》嘉靖十年条；《明史·兵志》及相关列传；边镇、地方志或对方材料 互相核对；不过，译写候选以编年材料定位；实录记载反映朝廷选择，崇祯期须另以奏疏、地方及敌对政权材料交叉。 因而这里标示的是材料能够支持的行动与制度位置，而不是对动机、地方执行或普通人经验的无保留复原。

\eventanchor{ML-1531-002}{明·1531（嘉靖十年）}
\paragraph{1531（嘉靖十年）：看到了一颗不吉利的彗星} 这一节点应放回本章已经展开的权力、财政、地方社会与边地关系中理解。账本把它出现的条件概括为：继承、官僚协调或皇权运作的压力使问题进入朝廷程序。 随后可追踪的变化是：它影响后续人事与政策议程；动机和派系标签不能由单一叙事裁定。 可由《世宗实录》嘉靖十年条；《明史·本纪》及相关列传；诏令、奏疏或会典版本 互相核对；不过，译写候选以编年材料定位；实录记载反映朝廷选择，崇祯期须另以奏疏、地方及敌对政权材料交叉。 因而这里标示的是材料能够支持的行动与制度位置，而不是对动机、地方执行或普通人经验的无保留复原。

\eventanchor{ML-1531-003}{明·1531（嘉靖十年）}
\paragraph{1531（嘉靖十年）：赋役折色、盐课、河工或田土核查的本年文书为财政改革史提供节点，制度名称不可跨区套用。} 这一节点应放回本章已经展开的权力、财政、地方社会与边地关系中理解。账本把它出现的条件概括为：财政征解、交通工程或货币支付的压力使相关措施进入政务。 随后可追踪的变化是：其法定安排不等于地方执行，后续须以地域材料追踪负担与成效。 可由《世宗实录》嘉靖十年条；《明会典》相关条目；地方志、黄册/契约/河工材料 互相核对；不过，桥接条目只标示可回查的年度问题与材料链；实录有中央选择性，崇祯期尤其需要奏疏、地方、朝鲜及满文材料交叉。 因而这里标示的是材料能够支持的行动与制度位置，而不是对动机、地方执行或普通人经验的无保留复原。

\eventanchor{ML-1531-004}{明·1531（嘉靖十年）}
\paragraph{1531（嘉靖十年）：嘉靖朝礼制、内阁与人事的本年诏令和奏议形成中枢协调线索，留中与施行之间应予区分。} 这一节点应放回本章已经展开的权力、财政、地方社会与边地关系中理解。账本把它出现的条件概括为：继承、官僚协调或皇权运作的压力使问题进入朝廷程序。 随后可追踪的变化是：它影响后续人事与政策议程；动机和派系标签不能由单一叙事裁定。 可由《世宗实录》嘉靖十年条；《明史·本纪》及相关列传；诏令、奏疏或会典版本 互相核对；不过，桥接条目只标示可回查的年度问题与材料链；实录有中央选择性，崇祯期尤其需要奏疏、地方、朝鲜及满文材料交叉。 因而这里标示的是材料能够支持的行动与制度位置，而不是对动机、地方执行或普通人经验的无保留复原。

\eventanchor{ML-1547-001}{明·1547（嘉靖二十六年）}
\paragraph{1547（嘉靖二十六年）：嘉靖倭口搜查：监察员报告称，东南沿海的海盗活动已失控。} 这一节点应放回本章已经展开的权力、财政、地方社会与边地关系中理解。账本把它出现的条件概括为：朝贡名义、边境安全与港口贸易的利益使交涉持续发生。 随后可追踪的变化是：它为后续册封、互市或海防安排留下文书与跨政权比较线索。 可由《世宗实录》嘉靖二十六年条；《明史·外国传》；《朝鲜王朝实录》、琉球《历代宝案》或海外档案 互相核对；不过，译写候选以编年材料定位；实录记载反映朝廷选择，崇祯期须另以奏疏、地方及敌对政权材料交叉。 因而这里标示的是材料能够支持的行动与制度位置，而不是对动机、地方执行或普通人经验的无保留复原。

\eventanchor{ML-1547-002}{明·1547（嘉靖二十六年）}
\paragraph{1547（嘉靖二十六年）：北边警报、边镇防务和西南处置的本年军政材料标记边疆压力，敌我称谓与军数应回到原文。} 这一节点应放回本章已经展开的权力、财政、地方社会与边地关系中理解。账本把它出现的条件概括为：前期边防、地方武装或跨区域资源调动的压力推动了该行动。 随后可追踪的变化是：它改变了后续调兵、补给或地方秩序的条件；战果和伤亡须分开核对。 可由《世宗实录》嘉靖二十六年条；《明史·兵志》及相关列传；边镇、地方志或对方材料 互相核对；不过，桥接条目只标示可回查的年度问题与材料链；实录有中央选择性，崇祯期尤其需要奏疏、地方、朝鲜及满文材料交叉。 因而这里标示的是材料能够支持的行动与制度位置，而不是对动机、地方执行或普通人经验的无保留复原。

\eventanchor{ML-1547-003}{明·1547（嘉靖二十六年）}
\paragraph{1547（嘉靖二十六年）：祀典、学校、出版或讲学的本年材料记录文化制度活动，规范话语不等于社会普遍认同。} 这一节点应放回本章已经展开的权力、财政、地方社会与边地关系中理解。账本把它出现的条件概括为：制度、教育或知识传播的需要推动了相关行动或文本形成。 随后可追踪的变化是：它提供制度和传播史的锚点，不能据此推断社会整体接受。 可由《世宗实录》嘉靖二十六年条；《明史·选举志》或《艺文志》；文集、碑刻或书目材料 互相核对；不过，桥接条目只标示可回查的年度问题与材料链；实录有中央选择性，崇祯期尤其需要奏疏、地方、朝鲜及满文材料交叉。 因而这里标示的是材料能够支持的行动与制度位置，而不是对动机、地方执行或普通人经验的无保留复原。

\eventanchor{MM-1547-ZHU-WAN-APPOINTMENT}{明·1547（嘉靖二十六年）}
\paragraph{1547（嘉靖二十六年）：朱纨出任巡抚浙江并兼理浙闽海防，以严禁私通海外和整饬港口为职} 这一节点应放回本章已经展开的权力、财政、地方社会与边地关系中理解。账本把它出现的条件概括为：双屿等港口的跨海贸易和武装保护网络扩张，地方官府既担忧海防又难以割断商人、渔民和中介的生计联系。 随后可追踪的变化是：朱纨的强硬政策加剧沿海利益冲突，并使海防指挥、地方官与中央言官之间的责任界限成为政治问题。 可由《世宗实录》嘉靖二十六年浙江条；《明史·朱纨传》；朱纨奏疏与《筹海图编》 互相核对；不过，任命及其法定职掌可靠；官府所谓私通、海寇和海商包含不同身份群体，禁令范围与沿海实际执行需以港口、商人和外文材料校验。 因而这里标示的是材料能够支持的行动与制度位置，而不是对动机、地方执行或普通人经验的无保留复原。

\eventanchor{ML-1548-001}{明·1548（嘉靖二十七年）二月}
\paragraph{1548（嘉靖二十七年）二月：嘉靖倭口劫掠：海盗劫掠宁波、台州。} 这一节点应放回本章已经展开的权力、财政、地方社会与边地关系中理解。账本把它出现的条件概括为：前期边防、地方武装或跨区域资源调动的压力推动了该行动。 随后可追踪的变化是：它改变了后续调兵、补给或地方秩序的条件；战果和伤亡须分开核对。 可由《世宗实录》嘉靖二十七年条；《明史·兵志》及相关列传；边镇、地方志或对方材料 互相核对；不过，译写候选以编年材料定位；实录记载反映朝廷选择，崇祯期须另以奏疏、地方及敌对政权材料交叉。 因而这里标示的是材料能够支持的行动与制度位置，而不是对动机、地方执行或普通人经验的无保留复原。

\eventanchor{ML-1548-002}{明·1548（嘉靖二十七年）四月}
\paragraph{1548（嘉靖二十七年）四月：嘉靖倭寇袭击：明军进攻双鱼，但港口内许多船只逃脱。} 这一节点应放回本章已经展开的权力、财政、地方社会与边地关系中理解。账本把它出现的条件概括为：前期边防、地方武装或跨区域资源调动的压力推动了该行动。 随后可追踪的变化是：它改变了后续调兵、补给或地方秩序的条件；战果和伤亡须分开核对。 可由《世宗实录》嘉靖二十七年条；《明史·兵志》及相关列传；边镇、地方志或对方材料 互相核对；不过，译写候选以编年材料定位；实录记载反映朝廷选择，崇祯期须另以奏疏、地方及敌对政权材料交叉。 因而这里标示的是材料能够支持的行动与制度位置，而不是对动机、地方执行或普通人经验的无保留复原。

\subsection{1548年前后的材料线索}

\eventanchor{ML-1548-003}{明·1548（嘉靖二十七年）六月}
\paragraph{1548（嘉靖二十七年）六月：蒙古人在宣府击败明军。} 这一节点应放回本章已经展开的权力、财政、地方社会与边地关系中理解。账本把它出现的条件概括为：前期边防、地方武装或跨区域资源调动的压力推动了该行动。 随后可追踪的变化是：它改变了后续调兵、补给或地方秩序的条件；战果和伤亡须分开核对。 可由《世宗实录》嘉靖二十七年条；《明史·兵志》及相关列传；边镇、地方志或对方材料 互相核对；不过，译写候选以编年材料定位；实录记载反映朝廷选择，崇祯期须另以奏疏、地方及敌对政权材料交叉。 因而这里标示的是材料能够支持的行动与制度位置，而不是对动机、地方执行或普通人经验的无保留复原。

\eventanchor{ML-1548-004}{明·1548（嘉靖二十七年）十月}
\paragraph{1548（嘉靖二十七年）十月：蒙古人袭击怀来。} 这一节点应放回本章已经展开的权力、财政、地方社会与边地关系中理解。账本把它出现的条件概括为：前期边防、地方武装或跨区域资源调动的压力推动了该行动。 随后可追踪的变化是：它改变了后续调兵、补给或地方秩序的条件；战果和伤亡须分开核对。 可由《世宗实录》嘉靖二十七年条；《明史·兵志》及相关列传；边镇、地方志或对方材料 互相核对；不过，译写候选以编年材料定位；实录记载反映朝廷选择，崇祯期须另以奏疏、地方及敌对政权材料交叉。 因而这里标示的是材料能够支持的行动与制度位置，而不是对动机、地方执行或普通人经验的无保留复原。

\eventanchor{ML-1548-005}{明·1548（嘉靖二十七年）}
\paragraph{1548（嘉靖二十七年）：明军开始配备火绳枪。} 这一节点应放回本章已经展开的权力、财政、地方社会与边地关系中理解。账本把它出现的条件概括为：继承、官僚协调或皇权运作的压力使问题进入朝廷程序。 随后可追踪的变化是：它影响后续人事与政策议程；动机和派系标签不能由单一叙事裁定。 可由《世宗实录》嘉靖二十七年条；《明史·本纪》及相关列传；诏令、奏疏或会典版本 互相核对；不过，译写候选以编年材料定位；实录记载反映朝廷选择，崇祯期须另以奏疏、地方及敌对政权材料交叉。 因而这里标示的是材料能够支持的行动与制度位置，而不是对动机、地方执行或普通人经验的无保留复原。

\eventanchor{MM-1548-XIA-YAN-EXECUTION}{明·1548（嘉靖二十七年）}
\paragraph{1548（嘉靖二十七年）：夏言因边务案件被处死} 这一节点应放回本章已经展开的权力、财政、地方社会与边地关系中理解。账本把它出现的条件概括为：围绕收复河套的边策争论失败，曾铣被杀后夏言失去政治支撑，严嵩得以发动攻讦。 随后可追踪的变化是：严嵩在内阁的优势进一步巩固；对北边主动战略的讨论受挫，朝臣更谨慎处理边务责任。 可由《世宗实录》嘉靖二十七年条；《明史·夏言传》；夏言、曾铣奏疏与相关题本 互相核对；不过，夏言被处死及案件程序可靠；罪名所涉军事责任和严嵩作用，应区别皇帝裁决、政敌攻讦与后世归因。 因而这里标示的是材料能够支持的行动与制度位置，而不是对动机、地方执行或普通人经验的无保留复原。

\eventanchor{ML-1549-001}{明·1549（嘉靖二十八年）}
\paragraph{1549（嘉靖二十八年）：赋役折色、盐课、河工或田土核查的本年文书为财政改革史提供节点，制度名称不可跨区套用。} 这一节点应放回本章已经展开的权力、财政、地方社会与边地关系中理解。账本把它出现的条件概括为：财政征解、交通工程或货币支付的压力使相关措施进入政务。 随后可追踪的变化是：其法定安排不等于地方执行，后续须以地域材料追踪负担与成效。 可由《世宗实录》嘉靖二十八年条；《明会典》相关条目；地方志、黄册/契约/河工材料 互相核对；不过，桥接条目只标示可回查的年度问题与材料链；实录有中央选择性，崇祯期尤其需要奏疏、地方、朝鲜及满文材料交叉。 因而这里标示的是材料能够支持的行动与制度位置，而不是对动机、地方执行或普通人经验的无保留复原。

\eventanchor{ML-1549-002}{明·1549（嘉靖二十八年）}
\paragraph{1549（嘉靖二十八年）：灾异、赈济及地方工程的本年报告可与州县材料互校，救济命令不自动证明救济到达。} 这一节点应放回本章已经展开的权力、财政、地方社会与边地关系中理解。账本把它出现的条件概括为：自然风险与地方报告促成朝廷或地方的应对记录。 随后可追踪的变化是：赈济、迁徙和损失的实际范围仍须以受灾地区材料分别核验。 可由《世宗实录》嘉靖二十八年条；《明史·五行志》；受灾府县地方志与奏疏 互相核对；不过，桥接条目只标示可回查的年度问题与材料链；实录有中央选择性，崇祯期尤其需要奏疏、地方、朝鲜及满文材料交叉。 因而这里标示的是材料能够支持的行动与制度位置，而不是对动机、地方执行或普通人经验的无保留复原。

\eventanchor{MM-1549-DATONG-MUTINY}{明·1549（嘉靖二十八年）}
\paragraph{1549（嘉靖二十八年）：大同军镇因军饷、军纪等问题发生兵变，朝廷调兵处置} 这一节点应放回本章已经展开的权力、财政、地方社会与边地关系中理解。账本把它出现的条件概括为：边镇军额虚耗、粮饷转运困难和将领—士卒关系紧张，使既有防务体系内部脆弱性暴露。 随后可追踪的变化是：朝廷以镇压和人事调整恢复秩序，但军饷与军纪问题未被根除，北边防务在翌年庚戌之变中再受检验。 可由《世宗实录》嘉靖二十八年大同条；《明史·兵志》；大同镇志与边镇奏疏 互相核对；不过，兵变与朝廷应对可靠；起事者的诉求和军饷亏欠数额须避免仅采纳镇将奏报或案卷定性。 因而这里标示的是材料能够支持的行动与制度位置，而不是对动机、地方执行或普通人经验的无保留复原。

\eventanchor{MM-1549-ZHU-WAN-DEATH}{明·1549（嘉靖二十八年）}
\paragraph{1549（嘉靖二十八年）：朱纨在遭弹劾、待讯的政治压力下自尽，严厉海禁执行路线受到挫折} 这一节点应放回本章已经展开的权力、财政、地方社会与边地关系中理解。账本把它出现的条件概括为：朱纨的港口查禁触及沿海商人、地方官和朝臣利益，处置俘获者的合法性又成为言官攻讦焦点。 随后可追踪的变化是：海防政策转向更依赖地方协商、招抚和后续督师安排；走私与海上武装网络并未因一名官员去世而消散。 可由《世宗实录》嘉靖二十八年条；《明史·朱纨传》；朱纨遗疏与相关题本 互相核对；不过，自尽及此前遭劾可由实录、传记和奏疏互证；具体指控、捕杀人数和自尽动机均来自高度政治化案卷，不能直接裁定。 因而这里标示的是材料能够支持的行动与制度位置，而不是对动机、地方执行或普通人经验的无保留复原。

\eventanchor{ML-1550-001}{明·1550（嘉靖二十九年）十月1日}
\paragraph{1550（嘉靖二十九年）十月1日：阿勒坦汗掠夺北京郊区} 这一节点应放回本章已经展开的权力、财政、地方社会与边地关系中理解。账本把它出现的条件概括为：继承、官僚协调或皇权运作的压力使问题进入朝廷程序。 随后可追踪的变化是：它影响后续人事与政策议程；动机和派系标签不能由单一叙事裁定。 可由《世宗实录》嘉靖二十九年条；《明史·本纪》及相关列传；诏令、奏疏或会典版本 互相核对；不过，译写候选以编年材料定位；实录记载反映朝廷选择，崇祯期须另以奏疏、地方及敌对政权材料交叉。 因而这里标示的是材料能够支持的行动与制度位置，而不是对动机、地方执行或普通人经验的无保留复原。

\eventanchor{ML-1550-002}{明·1550（嘉靖二十九年）十月6日}
\paragraph{1550（嘉靖二十九年）十月6日：明军被蒙古人击败} 这一节点应放回本章已经展开的权力、财政、地方社会与边地关系中理解。账本把它出现的条件概括为：前期边防、地方武装或跨区域资源调动的压力推动了该行动。 随后可追踪的变化是：它改变了后续调兵、补给或地方秩序的条件；战果和伤亡须分开核对。 可由《世宗实录》嘉靖二十九年条；《明史·兵志》及相关列传；边镇、地方志或对方材料 互相核对；不过，译写候选以编年材料定位；实录记载反映朝廷选择，崇祯期须另以奏疏、地方及敌对政权材料交叉。 因而这里标示的是材料能够支持的行动与制度位置，而不是对动机、地方执行或普通人经验的无保留复原。

\eventanchor{ML-1550-003}{明·1550（嘉靖二十九年）}
\paragraph{1550（嘉靖二十九年）：浙江村镇竖起栅栏应对土匪} 这一节点应放回本章已经展开的权力、财政、地方社会与边地关系中理解。账本把它出现的条件概括为：继承、官僚协调或皇权运作的压力使问题进入朝廷程序。 随后可追踪的变化是：它影响后续人事与政策议程；动机和派系标签不能由单一叙事裁定。 可由《世宗实录》嘉靖二十九年条；《明史·本纪》及相关列传；诏令、奏疏或会典版本 互相核对；不过，译写候选以编年材料定位；实录记载反映朝廷选择，崇祯期须另以奏疏、地方及敌对政权材料交叉。 因而这里标示的是材料能够支持的行动与制度位置，而不是对动机、地方执行或普通人经验的无保留复原。

\eventanchor{ML-1550-004}{明·1550（嘉靖二十九年）}
\paragraph{1550（嘉靖二十九年）：封贡、互市、月港和港口管理的本年文书记录边海关系重组，官方许可不等于贸易全貌。} 这一节点应放回本章已经展开的权力、财政、地方社会与边地关系中理解。账本把它出现的条件概括为：朝贡名义、边境安全与港口贸易的利益使交涉持续发生。 随后可追踪的变化是：它为后续册封、互市或海防安排留下文书与跨政权比较线索。 可由《世宗实录》嘉靖二十九年条；《明史·外国传》；《朝鲜王朝实录》、琉球《历代宝案》或海外档案 互相核对；不过，桥接条目只标示可回查的年度问题与材料链；实录有中央选择性，崇祯期尤其需要奏疏、地方、朝鲜及满文材料交叉。 因而这里标示的是材料能够支持的行动与制度位置，而不是对动机、地方执行或普通人经验的无保留复原。

\eventanchor{ML-1550-005}{明·1550（嘉靖二十九年）}
\paragraph{1550（嘉靖二十九年）：本年灾情、赈恤或河工记录连接中央与地方社会，区域差异需以府县材料追踪。} 这一节点应放回本章已经展开的权力、财政、地方社会与边地关系中理解。账本把它出现的条件概括为：自然风险与地方报告促成朝廷或地方的应对记录。 随后可追踪的变化是：赈济、迁徙和损失的实际范围仍须以受灾地区材料分别核验。 可由《世宗实录》嘉靖二十九年条；《明史·五行志》；受灾府县地方志与奏疏 互相核对；不过，桥接条目只标示可回查的年度问题与材料链；实录有中央选择性，崇祯期尤其需要奏疏、地方、朝鲜及满文材料交叉。 因而这里标示的是材料能够支持的行动与制度位置，而不是对动机、地方执行或普通人经验的无保留复原。

\eventanchor{MM-1550-GENGXU-RAID}{明·1550（嘉靖二十九年）八月}
\paragraph{1550（嘉靖二十九年）八月：俺答部抵达北京近郊，京师戒严，史称庚戌之变} 这一节点应放回本章已经展开的权力、财政、地方社会与边地关系中理解。账本把它出现的条件概括为：北边互市长期受阻，宣大防线又受军饷、指挥和补给问题制约，俺答部得以穿越防御空隙南下。 随后可追踪的变化是：京师防务、边墙修筑与边饷成为紧迫议题；明廷仍未解决封贡互市的结构性矛盾。 可由《世宗实录》嘉靖二十九年八月条；《明史·鞑靼传》；宣大边镇奏疏 互相核对；不过，抵京近郊和戒严可靠；双方兵数、破坏范围和谈判细节在边报与后出史书中差异显著。 因而这里标示的是材料能够支持的行动与制度位置，而不是对动机、地方执行或普通人经验的无保留复原。

\eventanchor{ML-1551-001}{明·1551（嘉靖三十年）}
\paragraph{1551（嘉靖三十年）：禁止渔船出海} 这一节点应放回本章已经展开的权力、财政、地方社会与边地关系中理解。账本把它出现的条件概括为：继承、官僚协调或皇权运作的压力使问题进入朝廷程序。 随后可追踪的变化是：它影响后续人事与政策议程；动机和派系标签不能由单一叙事裁定。 可由《世宗实录》嘉靖三十年条；《明史·本纪》及相关列传；诏令、奏疏或会典版本 互相核对；不过，译写候选以编年材料定位；实录记载反映朝廷选择，崇祯期须另以奏疏、地方及敌对政权材料交叉。 因而这里标示的是材料能够支持的行动与制度位置，而不是对动机、地方执行或普通人经验的无保留复原。

\eventanchor{ML-1551-002}{明·1551（嘉靖三十年）}
\paragraph{1551（嘉靖三十年）：严嵩时期至隆庆初的人事和奏议材料标记中枢权力变化，案狱与党争标签需保留来源立场。（材料核查重点：诏令或奏议的形成与发受链）} 这一节点应放回本章已经展开的权力、财政、地方社会与边地关系中理解。账本把它出现的条件概括为：继承、官僚协调或皇权运作的压力使问题进入朝廷程序。 随后可追踪的变化是：它影响后续人事与政策议程；动机和派系标签不能由单一叙事裁定。 可由《世宗实录》嘉靖三十年条；《明史·本纪》及相关列传；诏令、奏疏或会典版本 互相核对；不过，桥接条目只标示可回查的年度问题与材料链；实录有中央选择性，崇祯期尤其需要奏疏、地方、朝鲜及满文材料交叉。；本条特别核对诏令或奏议的形成与发受链。 因而这里标示的是材料能够支持的行动与制度位置，而不是对动机、地方执行或普通人经验的无保留复原。

\eventanchor{ML-1551-003}{明·1551（嘉靖三十年）}
\paragraph{1551（嘉靖三十年）：俺答边患、东南海防或沿海武装的本年军令与战报构成危机处理线索，战果和参与者须交叉核对。（材料核查重点：诏令或奏议的形成与发受链）} 这一节点应放回本章已经展开的权力、财政、地方社会与边地关系中理解。账本把它出现的条件概括为：前期边防、地方武装或跨区域资源调动的压力推动了该行动。 随后可追踪的变化是：它改变了后续调兵、补给或地方秩序的条件；战果和伤亡须分开核对。 可由《世宗实录》嘉靖三十年条；《明史·兵志》及相关列传；边镇、地方志或对方材料 互相核对；不过，桥接条目只标示可回查的年度问题与材料链；实录有中央选择性，崇祯期尤其需要奏疏、地方、朝鲜及满文材料交叉。；本条特别核对诏令或奏议的形成与发受链。 因而这里标示的是材料能够支持的行动与制度位置，而不是对动机、地方执行或普通人经验的无保留复原。

\eventanchor{ML-1551-004}{明·1551（嘉靖三十年）}
\paragraph{1551（嘉靖三十年）：严嵩时期至隆庆初的人事和奏议材料标记中枢权力变化，案狱与党争标签需保留来源立场。（材料核查重点：报称信息的来源、抄传与行政分类）} 这一节点应放回本章已经展开的权力、财政、地方社会与边地关系中理解。账本把它出现的条件概括为：继承、官僚协调或皇权运作的压力使问题进入朝廷程序。 随后可追踪的变化是：它影响后续人事与政策议程；动机和派系标签不能由单一叙事裁定。 可由《世宗实录》嘉靖三十年条；《明史·本纪》及相关列传；诏令、奏疏或会典版本 互相核对；不过，桥接条目只标示可回查的年度问题与材料链；实录有中央选择性，崇祯期尤其需要奏疏、地方、朝鲜及满文材料交叉。；本条特别核对报称信息的来源、抄传与行政分类。 因而这里标示的是材料能够支持的行动与制度位置，而不是对动机、地方执行或普通人经验的无保留复原。

\eventanchor{ML-1551-005}{明·1551（嘉靖三十年）}
\paragraph{1551（嘉靖三十年）：俺答边患、东南海防或沿海武装的本年军令与战报构成危机处理线索，战果和参与者须交叉核对。（材料核查重点：报称信息的来源、抄传与行政分类）} 这一节点应放回本章已经展开的权力、财政、地方社会与边地关系中理解。账本把它出现的条件概括为：前期边防、地方武装或跨区域资源调动的压力推动了该行动。 随后可追踪的变化是：它改变了后续调兵、补给或地方秩序的条件；战果和伤亡须分开核对。 可由《世宗实录》嘉靖三十年条；《明史·兵志》及相关列传；边镇、地方志或对方材料 互相核对；不过，桥接条目只标示可回查的年度问题与材料链；实录有中央选择性，崇祯期尤其需要奏疏、地方、朝鲜及满文材料交叉。；本条特别核对报称信息的来源、抄传与行政分类。 因而这里标示的是材料能够支持的行动与制度位置，而不是对动机、地方执行或普通人经验的无保留复原。

\eventanchor{MM-1551-ALTAN-SHANXI-RAID}{明·1551（嘉靖三十年）}
\paragraph{1551（嘉靖三十年）：俺答部再入山西北部活动，明廷调集边兵并讨论防线与互市} 这一节点应放回本章已经展开的权力、财政、地方社会与边地关系中理解。账本把它出现的条件概括为：1550年危机后封贡互市仍未解决，边镇饷运和防御纵深的缺陷使俺答部可继续以军事压力争取贸易条件。 随后可追踪的变化是：北部州县的防务与转运负担增加，朝廷更迫切讨论修边、增饷和市贸；冲突频率并未让双方停止接触。 可由《世宗实录》嘉靖三十年宣大与山西条；《明史·鞑靼传》；山西镇与大同镇志 互相核对；不过，边报、调兵和警报可证；入掠路线、损失和双方兵数多由边镇急报形成，不能与实际到兵、战果混为一谈。 因而这里标示的是材料能够支持的行动与制度位置，而不是对动机、地方执行或普通人经验的无保留复原。

\subsection{1552年前后的材料线索}

\eventanchor{ML-1552-001}{明·1552（嘉靖三十一年）四月}
\paragraph{1552（嘉靖三十一年）四月：明军在大同以北被蒙古人击败} 这一节点应放回本章已经展开的权力、财政、地方社会与边地关系中理解。账本把它出现的条件概括为：前期边防、地方武装或跨区域资源调动的压力推动了该行动。 随后可追踪的变化是：它改变了后续调兵、补给或地方秩序的条件；战果和伤亡须分开核对。 可由《世宗实录》嘉靖三十一年条；《明史·兵志》及相关列传；边镇、地方志或对方材料 互相核对；不过，译写候选以编年材料定位；实录记载反映朝廷选择，崇祯期须另以奏疏、地方及敌对政权材料交叉。 因而这里标示的是材料能够支持的行动与制度位置，而不是对动机、地方执行或普通人经验的无保留复原。

\eventanchor{ML-1552-002}{明·1552（嘉靖三十一年）}
\paragraph{1552（嘉靖三十一年）：嘉靖倭口突袭：突袭方攻击浙江沿海} 这一节点应放回本章已经展开的权力、财政、地方社会与边地关系中理解。账本把它出现的条件概括为：前期边防、地方武装或跨区域资源调动的压力推动了该行动。 随后可追踪的变化是：它改变了后续调兵、补给或地方秩序的条件；战果和伤亡须分开核对。 可由《世宗实录》嘉靖三十一年条；《明史·兵志》及相关列传；边镇、地方志或对方材料 互相核对；不过，译写候选以编年材料定位；实录记载反映朝廷选择，崇祯期须另以奏疏、地方及敌对政权材料交叉。 因而这里标示的是材料能够支持的行动与制度位置，而不是对动机、地方执行或普通人经验的无保留复原。

\eventanchor{ML-1552-003}{明·1552（嘉靖三十一年）}
\paragraph{1552（嘉靖三十一年）：嘉靖皇帝选拔800名8岁至14岁的少女入宫} 这一节点应放回本章已经展开的权力、财政、地方社会与边地关系中理解。账本把它出现的条件概括为：继承、官僚协调或皇权运作的压力使问题进入朝廷程序。 随后可追踪的变化是：它影响后续人事与政策议程；动机和派系标签不能由单一叙事裁定。 可由《世宗实录》嘉靖三十一年条；《明史·本纪》及相关列传；诏令、奏疏或会典版本 互相核对；不过，译写候选以编年材料定位；实录记载反映朝廷选择，崇祯期须另以奏疏、地方及敌对政权材料交叉。 因而这里标示的是材料能够支持的行动与制度位置，而不是对动机、地方执行或普通人经验的无保留复原。

\eventanchor{ML-1552-004}{明·1552（嘉靖三十一年）}
\paragraph{1552（嘉靖三十一年）：海防知识、学校、科举与礼制的本年文本提供制度和知识史线索，方案不等于实际配置。} 这一节点应放回本章已经展开的权力、财政、地方社会与边地关系中理解。账本把它出现的条件概括为：制度、教育或知识传播的需要推动了相关行动或文本形成。 随后可追踪的变化是：它提供制度和传播史的锚点，不能据此推断社会整体接受。 可由《世宗实录》嘉靖三十一年条；《明史·选举志》或《艺文志》；文集、碑刻或书目材料 互相核对；不过，桥接条目只标示可回查的年度问题与材料链；实录有中央选择性，崇祯期尤其需要奏疏、地方、朝鲜及满文材料交叉。 因而这里标示的是材料能够支持的行动与制度位置，而不是对动机、地方执行或普通人经验的无保留复原。

\eventanchor{MM-1552-ALTAN-DATONG-RAID}{明·1552（嘉靖三十一年）}
\paragraph{1552（嘉靖三十一年）：俺答部在大同一线继续入掠，边镇以调兵、守城和转运应对} 这一节点应放回本章已经展开的权力、财政、地方社会与边地关系中理解。账本把它出现的条件概括为：前一年边患未获制度性解决，俺答部可沿熟悉的草原—关隘通道行动，大同军镇则受兵饷和指挥协调限制。 随后可追踪的变化是：大同的守备、粮饷和城堡维修再次成为中央议题；明蒙关系仍在军事冲突与互市谈判之间摆动。 可由《世宗实录》嘉靖三十一年大同条；《明史·鞑靼传》；大同镇志与边饷奏疏 互相核对；不过，战事警报和明廷应对可靠；边镇奏报中的破坏面积、俘斩和敌军数量具有报功或求援动机，须同蒙古材料及地方志核对。 因而这里标示的是材料能够支持的行动与制度位置，而不是对动机、地方执行或普通人经验的无保留复原。

\eventanchor{ML-1553-001}{明·1553（嘉靖三十二年）}
\paragraph{1553（嘉靖三十二年）：嘉靖倭寇劫掠：王直劫掠台州以北的浙江沿海} 这一节点应放回本章已经展开的权力、财政、地方社会与边地关系中理解。账本把它出现的条件概括为：前期边防、地方武装或跨区域资源调动的压力推动了该行动。 随后可追踪的变化是：它改变了后续调兵、补给或地方秩序的条件；战果和伤亡须分开核对。 可由《世宗实录》嘉靖三十二年条；《明史·兵志》及相关列传；边镇、地方志或对方材料 互相核对；不过，译写候选以编年材料定位；实录记载反映朝廷选择，崇祯期须另以奏疏、地方及敌对政权材料交叉。 因而这里标示的是材料能够支持的行动与制度位置，而不是对动机、地方执行或普通人经验的无保留复原。

\eventanchor{ML-1553-002}{明·1553（嘉靖三十二年）}
\paragraph{1553（嘉靖三十二年）：海防知识、学校、科举与礼制的本年文本提供制度和知识史线索，方案不等于实际配置。} 这一节点应放回本章已经展开的权力、财政、地方社会与边地关系中理解。账本把它出现的条件概括为：制度、教育或知识传播的需要推动了相关行动或文本形成。 随后可追踪的变化是：它提供制度和传播史的锚点，不能据此推断社会整体接受。 可由《世宗实录》嘉靖三十二年条；《明史·选举志》或《艺文志》；文集、碑刻或书目材料 互相核对；不过，桥接条目只标示可回查的年度问题与材料链；实录有中央选择性，崇祯期尤其需要奏疏、地方、朝鲜及满文材料交叉。 因而这里标示的是材料能够支持的行动与制度位置，而不是对动机、地方执行或普通人经验的无保留复原。

\eventanchor{ML-1553-003}{明·1553（嘉靖三十二年）}
\paragraph{1553（嘉靖三十二年）：军饷、盐课、漕运和赋役的本年行政材料显示中晚明财政压力，法定额与实收不可互代。} 这一节点应放回本章已经展开的权力、财政、地方社会与边地关系中理解。账本把它出现的条件概括为：财政征解、交通工程或货币支付的压力使相关措施进入政务。 随后可追踪的变化是：其法定安排不等于地方执行，后续须以地域材料追踪负担与成效。 可由《世宗实录》嘉靖三十二年条；《明会典》相关条目；地方志、黄册/契约/河工材料 互相核对；不过，桥接条目只标示可回查的年度问题与材料链；实录有中央选择性，崇祯期尤其需要奏疏、地方、朝鲜及满文材料交叉。 因而这里标示的是材料能够支持的行动与制度位置，而不是对动机、地方执行或普通人经验的无保留复原。

\eventanchor{ML-1553-004}{明·1553（嘉靖三十二年）}
\paragraph{1553（嘉靖三十二年）：本年灾情、赈恤或河工记录连接中央与地方社会，区域差异需以府县材料追踪。} 这一节点应放回本章已经展开的权力、财政、地方社会与边地关系中理解。账本把它出现的条件概括为：自然风险与地方报告促成朝廷或地方的应对记录。 随后可追踪的变化是：赈济、迁徙和损失的实际范围仍须以受灾地区材料分别核验。 可由《世宗实录》嘉靖三十二年条；《明史·五行志》；受灾府县地方志与奏疏 互相核对；不过，桥接条目只标示可回查的年度问题与材料链；实录有中央选择性，崇祯期尤其需要奏疏、地方、朝鲜及满文材料交叉。 因而这里标示的是材料能够支持的行动与制度位置，而不是对动机、地方执行或普通人经验的无保留复原。

\eventanchor{MM-1553-SHANGHAI-WALL}{明·1553（嘉靖三十二年）}
\paragraph{1553（嘉靖三十二年）：上海县为应对沿海袭扰修筑城墙，官府、士绅与居民参与筹办} 这一节点应放回本章已经展开的权力、财政、地方社会与边地关系中理解。账本把它出现的条件概括为：浙东、江南沿海武装活动威胁市镇和交通，地方社会希望取得可见的防御设施。 随后可追踪的变化是：城墙重塑县城空间与地方动员方式；它只能保护有限城区，不能单独解决沿海与内河的安全问题。 可由《上海县志》城池条；上海古城墙遗址调查资料；《明史·兵志》倭患相关条 互相核对；不过，筑城时间和工程可由方志、碑记互证；出资比例、劳役负担与城墙防御效能须以具体材料而非纪功文概括。 因而这里标示的是材料能够支持的行动与制度位置，而不是对动机、地方执行或普通人经验的无保留复原。

\eventanchor{ML-1554-001}{明·1554（嘉靖三十三年）春}
\paragraph{1554（嘉靖三十三年）春：北京爆发严重疫情} 这一节点应放回本章已经展开的权力、财政、地方社会与边地关系中理解。账本把它出现的条件概括为：自然风险与地方报告促成朝廷或地方的应对记录。 随后可追踪的变化是：赈济、迁徙和损失的实际范围仍须以受灾地区材料分别核验。 可由《世宗实录》嘉靖三十三年条；《明史·五行志》；受灾府县地方志与奏疏 互相核对；不过，译写候选以编年材料定位；实录记载反映朝廷选择，崇祯期须另以奏疏、地方及敌对政权材料交叉。 因而这里标示的是材料能够支持的行动与制度位置，而不是对动机、地方执行或普通人经验的无保留复原。

\eventanchor{ML-1554-002}{明·1554（嘉靖三十三年）三月}
\paragraph{1554（嘉靖三十三年）三月：嘉靖倭口突袭：海盗杀死松江知县，占领崇明岛} 这一节点应放回本章已经展开的权力、财政、地方社会与边地关系中理解。账本把它出现的条件概括为：前期边防、地方武装或跨区域资源调动的压力推动了该行动。 随后可追踪的变化是：它改变了后续调兵、补给或地方秩序的条件；战果和伤亡须分开核对。 可由《世宗实录》嘉靖三十三年条；《明史·兵志》及相关列传；边镇、地方志或对方材料 互相核对；不过，译写候选以编年材料定位；实录记载反映朝廷选择，崇祯期须另以奏疏、地方及敌对政权材料交叉。 因而这里标示的是材料能够支持的行动与制度位置，而不是对动机、地方执行或普通人经验的无保留复原。

\eventanchor{ML-1554-003}{明·1554（嘉靖三十三年）}
\paragraph{1554（嘉靖三十三年）：葡中协定（1554年）：莱昂内尔·德索萨贿赂海防副专员，让葡萄牙人留在澳门，每年支付500两银子，并对一半产品征收20\%的关税} 这一节点应放回本章已经展开的权力、财政、地方社会与边地关系中理解。账本把它出现的条件概括为：朝贡名义、边境安全与港口贸易的利益使交涉持续发生。 随后可追踪的变化是：它为后续册封、互市或海防安排留下文书与跨政权比较线索。 可由《世宗实录》嘉靖三十三年条；《明史·外国传》；《朝鲜王朝实录》、琉球《历代宝案》或海外档案 互相核对；不过，译写候选以编年材料定位；实录记载反映朝廷选择，崇祯期须另以奏疏、地方及敌对政权材料交叉。 因而这里标示的是材料能够支持的行动与制度位置，而不是对动机、地方执行或普通人经验的无保留复原。

\eventanchor{ML-1554-004}{明·1554（嘉靖三十三年）}
\paragraph{1554（嘉靖三十三年）：俺答边患、东南海防或沿海武装的本年军令与战报构成危机处理线索，战果和参与者须交叉核对。} 这一节点应放回本章已经展开的权力、财政、地方社会与边地关系中理解。账本把它出现的条件概括为：前期边防、地方武装或跨区域资源调动的压力推动了该行动。 随后可追踪的变化是：它改变了后续调兵、补给或地方秩序的条件；战果和伤亡须分开核对。 可由《世宗实录》嘉靖三十三年条；《明史·兵志》及相关列传；边镇、地方志或对方材料 互相核对；不过，桥接条目只标示可回查的年度问题与材料链；实录有中央选择性，崇祯期尤其需要奏疏、地方、朝鲜及满文材料交叉。 因而这里标示的是材料能够支持的行动与制度位置，而不是对动机、地方执行或普通人经验的无保留复原。

\eventanchor{MM-1554-WANGJIANGJING}{明·1554（嘉靖三十三年）}
\paragraph{1554（嘉靖三十三年）：张经、俞大猷等在王江泾击败一支海上武装} 这一节点应放回本章已经展开的权力、财政、地方社会与边地关系中理解。账本把它出现的条件概括为：海禁、私人贸易和沿海防务失调使武装集团可依托港汊、岛屿和市镇网络活动。 随后可追踪的变化是：胜利暂时缓解苏松浙北压力，但张经随后因政治案件被杀，海防指挥的连续性受损。 可由《世宗实录》嘉靖三十三年浙江条；《明史·张经传》；《筹海图编》与嘉兴地方志 互相核对；不过，战事和主要将领可靠；“倭寇”内部包含日本人、中国海商与沿海武装等不同成分，敌我人数和斩获不能照录捷报。 因而这里标示的是材料能够支持的行动与制度位置，而不是对动机、地方执行或普通人经验的无保留复原。

\eventanchor{ML-1555-001}{明·1555（嘉靖三十四年）}
\paragraph{1555（嘉靖三十四年）：嘉靖倭口劫掠：海盗袭击杭州} 这一节点应放回本章已经展开的权力、财政、地方社会与边地关系中理解。账本把它出现的条件概括为：前期边防、地方武装或跨区域资源调动的压力推动了该行动。 随后可追踪的变化是：它改变了后续调兵、补给或地方秩序的条件；战果和伤亡须分开核对。 可由《世宗实录》嘉靖三十四年条；《明史·兵志》及相关列传；边镇、地方志或对方材料 互相核对；不过，译写候选以编年材料定位；实录记载反映朝廷选择，崇祯期须另以奏疏、地方及敌对政权材料交叉。 因而这里标示的是材料能够支持的行动与制度位置，而不是对动机、地方执行或普通人经验的无保留复原。

\eventanchor{ML-1555-002}{明·1555（嘉靖三十四年）五月}
\paragraph{1555（嘉靖三十四年）五月：嘉靖倭寇突袭：明军击败嘉兴以北的一支大型突袭部队} 这一节点应放回本章已经展开的权力、财政、地方社会与边地关系中理解。账本把它出现的条件概括为：前期边防、地方武装或跨区域资源调动的压力推动了该行动。 随后可追踪的变化是：它改变了后续调兵、补给或地方秩序的条件；战果和伤亡须分开核对。 可由《世宗实录》嘉靖三十四年条；《明史·兵志》及相关列传；边镇、地方志或对方材料 互相核对；不过，译写候选以编年材料定位；实录记载反映朝廷选择，崇祯期须另以奏疏、地方及敌对政权材料交叉。 因而这里标示的是材料能够支持的行动与制度位置，而不是对动机、地方执行或普通人经验的无保留复原。

\eventanchor{ML-1555-003}{明·1555（嘉靖三十四年）}
\paragraph{1555（嘉靖三十四年）：嘉靖皇帝选拔180名10岁以下的少女入宫} 这一节点应放回本章已经展开的权力、财政、地方社会与边地关系中理解。账本把它出现的条件概括为：继承、官僚协调或皇权运作的压力使问题进入朝廷程序。 随后可追踪的变化是：它影响后续人事与政策议程；动机和派系标签不能由单一叙事裁定。 可由《世宗实录》嘉靖三十四年条；《明史·本纪》及相关列传；诏令、奏疏或会典版本 互相核对；不过，译写候选以编年材料定位；实录记载反映朝廷选择，崇祯期须另以奏疏、地方及敌对政权材料交叉。 因而这里标示的是材料能够支持的行动与制度位置，而不是对动机、地方执行或普通人经验的无保留复原。

\eventanchor{MM-1555-YANG-JISHENG-EXECUTION}{明·1555（嘉靖三十四年）}
\paragraph{1555（嘉靖三十四年）：杨继盛在狱中被处死，其弹劾严嵩的奏疏随后广泛流传} 这一节点应放回本章已经展开的权力、财政、地方社会与边地关系中理解。账本把它出现的条件概括为：杨继盛曾以边务和严嵩问题上疏，触犯严嵩集团并在司法程序中长期羁押。 随后可追踪的变化是：案件显示批评内阁权臣的高风险；遗稿和祠祀使其成为后世士大夫政治记忆的重要符号。 可由《世宗实录》嘉靖三十四年条；《明史·杨继盛传》；《杨忠愍公集》 互相核对；不过，杨继盛被杀及其奏疏文本可靠；将其单独解释为“清流”与“奸臣”对抗，会遮蔽边策、司法和嘉靖政治的具体脉络。 因而这里标示的是材料能够支持的行动与制度位置，而不是对动机、地方执行或普通人经验的无保留复原。

\eventanchor{MM-1555-ZHANG-JING-EXECUTION}{明·1555（嘉靖三十四年）}
\paragraph{1555（嘉靖三十四年）：张经在王江泾获胜后仍因海防案件被处死，东南督师人事骤变} 这一节点应放回本章已经展开的权力、财政、地方社会与边地关系中理解。账本把它出现的条件概括为：海防战事需要跨省调兵与等待战机，张经的作战节奏同中央急于求效和内阁竞争发生冲突。 随后可追踪的变化是：经验丰富的统帅被撤除，东南军事指挥转由胡宗宪等重新组织；战场胜负与朝廷问责之间的张力更加明显。 可由《世宗实录》嘉靖三十四年条；《明史·张经传》；张经、赵文华相关奏疏 互相核对；不过，张经被杀及官方罪名可靠；其是否延误军机、赵文华和严嵩所起作用以及战功归属，都须区分题本、案狱与后出的传记评价。 因而这里标示的是材料能够支持的行动与制度位置，而不是对动机、地方执行或普通人经验的无保留复原。

\subsection{1556年前后的材料线索}

\eventanchor{ML-1556-001}{明·1556（嘉靖三十五年）一月}
\paragraph{1556（嘉靖三十五年）一月：1556年陕西大地震：陕西发生地震，死亡人数超过80万} 这一节点应放回本章已经展开的权力、财政、地方社会与边地关系中理解。账本把它出现的条件概括为：自然风险与地方报告促成朝廷或地方的应对记录。 随后可追踪的变化是：赈济、迁徙和损失的实际范围仍须以受灾地区材料分别核验。 可由《世宗实录》嘉靖三十五年条；《明史·五行志》；受灾府县地方志与奏疏 互相核对；不过，译写候选以编年材料定位；实录记载反映朝廷选择，崇祯期须另以奏疏、地方及敌对政权材料交叉。 因而这里标示的是材料能够支持的行动与制度位置，而不是对动机、地方执行或普通人经验的无保留复原。

\eventanchor{ML-1556-002}{明·1556（嘉靖三十五年）}
\paragraph{1556（嘉靖三十五年）：嘉靖倭口劫掠：海盗袭击从南京到杭州的整个海岸线} 这一节点应放回本章已经展开的权力、财政、地方社会与边地关系中理解。账本把它出现的条件概括为：前期边防、地方武装或跨区域资源调动的压力推动了该行动。 随后可追踪的变化是：它改变了后续调兵、补给或地方秩序的条件；战果和伤亡须分开核对。 可由《世宗实录》嘉靖三十五年条；《明史·兵志》及相关列传；边镇、地方志或对方材料 互相核对；不过，译写候选以编年材料定位；实录记载反映朝廷选择，崇祯期须另以奏疏、地方及敌对政权材料交叉。 因而这里标示的是材料能够支持的行动与制度位置，而不是对动机、地方执行或普通人经验的无保留复原。

\eventanchor{ML-1556-003}{明·1556（嘉靖三十五年）}
\paragraph{1556（嘉靖三十五年）：嘉靖皇帝请求礼部寻找一些神奇的植物来让他长生不老} 这一节点应放回本章已经展开的权力、财政、地方社会与边地关系中理解。账本把它出现的条件概括为：继承、官僚协调或皇权运作的压力使问题进入朝廷程序。 随后可追踪的变化是：它影响后续人事与政策议程；动机和派系标签不能由单一叙事裁定。 可由《世宗实录》嘉靖三十五年条；《明史·本纪》及相关列传；诏令、奏疏或会典版本 互相核对；不过，译写候选以编年材料定位；实录记载反映朝廷选择，崇祯期须另以奏疏、地方及敌对政权材料交叉。 因而这里标示的是材料能够支持的行动与制度位置，而不是对动机、地方执行或普通人经验的无保留复原。

\eventanchor{ML-1556-004}{明·1556（嘉靖三十五年）}
\paragraph{1556（嘉靖三十五年）：军饷、盐课、漕运和赋役的本年行政材料显示中晚明财政压力，法定额与实收不可互代。} 这一节点应放回本章已经展开的权力、财政、地方社会与边地关系中理解。账本把它出现的条件概括为：财政征解、交通工程或货币支付的压力使相关措施进入政务。 随后可追踪的变化是：其法定安排不等于地方执行，后续须以地域材料追踪负担与成效。 可由《世宗实录》嘉靖三十五年条；《明会典》相关条目；地方志、黄册/契约/河工材料 互相核对；不过，桥接条目只标示可回查的年度问题与材料链；实录有中央选择性，崇祯期尤其需要奏疏、地方、朝鲜及满文材料交叉。 因而这里标示的是材料能够支持的行动与制度位置，而不是对动机、地方执行或普通人经验的无保留复原。

\eventanchor{MM-1556-SHAANXI-EARTHQUAKE}{明·1556（嘉靖三十四年十二月）}
\paragraph{1556（嘉靖三十四年十二月）：华县大地震波及关中及邻近地区，朝廷下诏赈恤、修复} 这一节点应放回本章已经展开的权力、财政、地方社会与边地关系中理解。账本把它出现的条件概括为：关中人口集中于黄土高原窑居和城镇，强震使住房、交通和地方行政同时受损。 随后可追踪的变化是：赈恤和修复成为跨省财政、粮运和地方组织的长期任务；灾害影响不宜压缩为单一全国死亡数字。 可由《世宗实录》嘉靖三十四年十二月条；《明史·五行志》；关中地方志与地震碑刻 互相核对；不过，地震日期和广泛破坏可靠；死亡总数多为后世汇算，地方损失与赈济到达情况必须分别以方志、碑刻和文书核验。 因而这里标示的是材料能够支持的行动与制度位置，而不是对动机、地方执行或普通人经验的无保留复原。

\eventanchor{ML-1557-001}{明·1557（嘉靖三十六年）五月}
\paragraph{1557（嘉靖三十六年）五月：故宫三大殿毁于火灾} 这一节点应放回本章已经展开的权力、财政、地方社会与边地关系中理解。账本把它出现的条件概括为：自然风险与地方报告促成朝廷或地方的应对记录。 随后可追踪的变化是：赈济、迁徙和损失的实际范围仍须以受灾地区材料分别核验。 可由《世宗实录》嘉靖三十六年条；《明史·五行志》；受灾府县地方志与奏疏 互相核对；不过，译写候选以编年材料定位；实录记载反映朝廷选择，崇祯期须另以奏疏、地方及敌对政权材料交叉。 因而这里标示的是材料能够支持的行动与制度位置，而不是对动机、地方执行或普通人经验的无保留复原。

\eventanchor{ML-1557-002}{明·1557（嘉靖三十六年）冬}
\paragraph{1557（嘉靖三十六年）冬：阿勒坦汗之子僧格围攻大同附近的驻军} 这一节点应放回本章已经展开的权力、财政、地方社会与边地关系中理解。账本把它出现的条件概括为：前期边防、地方武装或跨区域资源调动的压力推动了该行动。 随后可追踪的变化是：它改变了后续调兵、补给或地方秩序的条件；战果和伤亡须分开核对。 可由《世宗实录》嘉靖三十六年条；《明史·兵志》及相关列传；边镇、地方志或对方材料 互相核对；不过，译写候选以编年材料定位；实录记载反映朝廷选择，崇祯期须另以奏疏、地方及敌对政权材料交叉。 因而这里标示的是材料能够支持的行动与制度位置，而不是对动机、地方执行或普通人经验的无保留复原。

\eventanchor{ML-1557-003}{明·1557（嘉靖三十六年）}
\paragraph{1557（嘉靖三十六年）：封贡、互市、月港和港口管理的本年文书记录边海关系重组，官方许可不等于贸易全貌。} 这一节点应放回本章已经展开的权力、财政、地方社会与边地关系中理解。账本把它出现的条件概括为：朝贡名义、边境安全与港口贸易的利益使交涉持续发生。 随后可追踪的变化是：它为后续册封、互市或海防安排留下文书与跨政权比较线索。 可由《世宗实录》嘉靖三十六年条；《明史·外国传》；《朝鲜王朝实录》、琉球《历代宝案》或海外档案 互相核对；不过，桥接条目只标示可回查的年度问题与材料链；实录有中央选择性，崇祯期尤其需要奏疏、地方、朝鲜及满文材料交叉。 因而这里标示的是材料能够支持的行动与制度位置，而不是对动机、地方执行或普通人经验的无保留复原。

\eventanchor{MM-1557-MACAU-SETTLEMENT}{明·1557（嘉靖三十六年）}
\paragraph{1557（嘉靖三十六年）：葡萄牙商人获准在澳门一带停泊、居留和贸易，地方官府对其征收地租等费用} 这一节点应放回本章已经展开的权力、财政、地方社会与边地关系中理解。账本把它出现的条件概括为：珠江口贸易需求和葡萄牙航海网络持续存在，广东地方官需在海防、税收和港口秩序之间寻求可操作安排。 随后可追踪的变化是：澳门逐渐成为中外商人、教士和转口贸易的节点，但其法律地位和明朝主权表述长期具有多重层次。 可由香山、澳门地方志早期条目；《明史·佛郎机传》；葡萄牙澳门档案与港口考古 互相核对；不过，澳门早期居留的形成大体可靠；“1557租地”是后出的简化说法，准入、纳税与实际控制权须以中葡档案分别讨论。 因而这里标示的是材料能够支持的行动与制度位置，而不是对动机、地方执行或普通人经验的无保留复原。

\eventanchor{MM-1557-WANG-ZHI-SURRENDER}{明·1557（嘉靖三十六年）}
\paragraph{1557（嘉靖三十六年）：王直在胡宗宪招抚安排下抵浙受审，跨海贸易网络中的一支核心集团被纳入官府处置} 这一节点应放回本章已经展开的权力、财政、地方社会与边地关系中理解。账本把它出现的条件概括为：连续海战难以截断岛屿、港口和海外市场之间的人员流动，胡宗宪试图以招抚分化王直集团并取得情报。 随后可追踪的变化是：王直被拘使招抚转入审讯和问责程序，部分网络继续运作；1559年处决不能被解释为沿海私贸的终结。 可由《世宗实录》嘉靖三十六年浙江条；《明史·王直传》；胡宗宪奏疏与葡日贸易材料 互相核对；不过，抵浙、受审和招抚过程可由实录、传记及海防书互见；王直是否自愿降附、随从人数和谈判承诺在官府与海外材料中表述不一。 因而这里标示的是材料能够支持的行动与制度位置，而不是对动机、地方执行或普通人经验的无保留复原。

\eventanchor{ML-1558-001}{明·1558（嘉靖三十七年）四月}
\paragraph{1558（嘉靖三十七年）四月：嘉靖倭口劫掠：海盗袭击浙江、闽北} 这一节点应放回本章已经展开的权力、财政、地方社会与边地关系中理解。账本把它出现的条件概括为：前期边防、地方武装或跨区域资源调动的压力推动了该行动。 随后可追踪的变化是：它改变了后续调兵、补给或地方秩序的条件；战果和伤亡须分开核对。 可由《世宗实录》嘉靖三十七年条；《明史·兵志》及相关列传；边镇、地方志或对方材料 互相核对；不过，译写候选以编年材料定位；实录记载反映朝廷选择，崇祯期须另以奏疏、地方及敌对政权材料交叉。 因而这里标示的是材料能够支持的行动与制度位置，而不是对动机、地方执行或普通人经验的无保留复原。

\eventanchor{ML-1558-002}{明·1558（嘉靖三十七年）}
\paragraph{1558（嘉靖三十七年）：援军到来后，僧格撤退} 这一节点应放回本章已经展开的权力、财政、地方社会与边地关系中理解。账本把它出现的条件概括为：前期边防、地方武装或跨区域资源调动的压力推动了该行动。 随后可追踪的变化是：它改变了后续调兵、补给或地方秩序的条件；战果和伤亡须分开核对。 可由《世宗实录》嘉靖三十七年条；《明史·兵志》及相关列传；边镇、地方志或对方材料 互相核对；不过，译写候选以编年材料定位；实录记载反映朝廷选择，崇祯期须另以奏疏、地方及敌对政权材料交叉。 因而这里标示的是材料能够支持的行动与制度位置，而不是对动机、地方执行或普通人经验的无保留复原。

\eventanchor{ML-1558-003}{明·1558（嘉靖三十七年）}
\paragraph{1558（嘉靖三十七年）：礼部向嘉靖皇帝进贡1860株神奇植物} 这一节点应放回本章已经展开的权力、财政、地方社会与边地关系中理解。账本把它出现的条件概括为：继承、官僚协调或皇权运作的压力使问题进入朝廷程序。 随后可追踪的变化是：它影响后续人事与政策议程；动机和派系标签不能由单一叙事裁定。 可由《世宗实录》嘉靖三十七年条；《明史·本纪》及相关列传；诏令、奏疏或会典版本 互相核对；不过，译写候选以编年材料定位；实录记载反映朝廷选择，崇祯期须另以奏疏、地方及敌对政权材料交叉。 因而这里标示的是材料能够支持的行动与制度位置，而不是对动机、地方执行或普通人经验的无保留复原。
